\documentclass[11pt,a4paper]{article}

\usepackage[utf8]{inputenc}
\usepackage[T1]{fontenc}
\usepackage[provide=*,english]{babel}
\usepackage{lmodern}
\usepackage{microtype}
\usepackage{amsmath,amssymb,mathtools}
\usepackage{geometry}
\usepackage{xcolor}
\usepackage{enumitem}
\usepackage{booktabs}
\usepackage{fancyhdr}
\usepackage[hidelinks]{hyperref}

\geometry{left=2.3cm,right=2.3cm,top=2.2cm,bottom=2.25cm,headheight=14pt}
\setlength{\parindent}{0pt}
\setlength{\parskip}{0.42em}
\setlist[itemize]{topsep=0.25em,itemsep=0.18em,leftmargin=1.5em}
\setlist[enumerate]{topsep=0.3em,itemsep=0.25em,leftmargin=1.9em}
\allowdisplaybreaks

\definecolor{blue}{RGB}{34,73,110}
\definecolor{gray}{RGB}{90,96,102}

\pagestyle{fancy}
\fancyhf{}
\fancyhead[L]{\small\color{gray}GC2D: ABBA4 implicit diagnostics}
\fancyhead[R]{\small\color{gray}Trajectory and tangent dynamics}
\fancyfoot[C]{\thepage}
\renewcommand{\headrulewidth}{0.3pt}

\newcommand{\R}{\mathbb{R}}
\newcommand{\Id}{I}
\newcommand{\norm}[1]{\left\lVert #1\right\rVert}

\begin{document}

\begin{center}
  {\LARGE\bfseries\color{blue}Jacobian Diagnostics for ABBA4 Implicit}\par
  \vspace{0.35em}
  {\large Three ideal-root tangents along one signed triple jump}\par
  \vspace{0.45em}
  {\small Physical GC2D trajectory, local symplecticity, and accumulated sensitivity}
\end{center}

\vspace{0.5em}

\section{Purpose of the diagnostics}

An \texttt{ABBA4Implicit} integration can be observed through the physical
state and an infinitesimal perturbation,
\[
 z_n\in\R^2,
 \qquad
 \xi_n\in\R^2.
\]
The trajectory advances with the complete three-solve method,
\[
 z_{n+1}=\Psi^{[4]}_{h,t_n}(z_n),
\]
while the tangent follows the derivative of that same nonlinear map,
\[
 \boxed{
 \xi_{n+1}=D\Psi^{[4]}_{h,t_n}(z_n)\,\xi_n.
 }
\]

The Jacobian is not a replacement for the state update.  In general,
\[
 z_{n+1}\neq D\Psi^{[4]}_{h,t_n}(z_n)z_n.
\]
It acts on a perturbation or on an accumulated derivative with respect to the
initial condition.

The production observer exposes the complete accepted outer step and an
ordered tuple of three signed substep snapshots.  The diagnostic first
constructs one ideal-root physical Jacobian from each snapshot and then applies
the chain rule.

\section{Signed substep geometry}

Let
\[
 (c_1,c_2,c_3)=(\gamma,\delta,\gamma),
\]
where
\[
 \gamma=\frac{1}{2-2^{1/3}},
 \qquad
 \delta=-\frac{2^{1/3}}{2-2^{1/3}}.
\]
Define
\[
 \ell_j=c_jh,
 \qquad
 \tau_j=t_n+h\sum_{r<j}c_r,
 \qquad
 s_j=\frac{\ell_j}{2}.
\]
Then
\[
 z^{[0]}=z_n,
 \qquad
 z^{[j]}=\Psi^{[2]}_{\ell_j,\tau_j}(z^{[j-1]}),
 \qquad
 z_{n+1}=z^{[3]}.
\]
The signed clock is
\begin{center}
\begin{tabular}{@{}cccc@{}}
\toprule
Substep & Start & Duration & End\\
\midrule
$1$ & $t_n$ & $\gamma h$ & $t_n+\gamma h$\\
$2$ & $t_n+\gamma h$ & $\delta h$ & $t_n+(\gamma+\delta)h$\\
$3$ & $t_n+(\gamma+\delta)h$ & $\gamma h$ & $t_n+h$\\
\bottomrule
\end{tabular}
\end{center}
The second row runs backward in time.  A diagnostic must use the stored signed
duration and start time rather than infer monotone stage times from the outer
interval.

\section{One projected ABBA substep}

For a provisional multiplier $\mu_j\in\R^2$, initialize
\[
 u_j^{(0)}=z^{[j-1]}+\mu_j,
 \qquad
 v_j^{(0)}=z^{[j-1]}-\mu_j.
\]
The endpoint-time stages are
\begin{align*}
 u_j^{(1)}&=u_j^{(0)}+s_jf(v_j^{(0)},\tau_j),\\
 v_j^{(1)}&=v_j^{(0)}+s_jf(u_j^{(1)},\tau_j),\\
 v_j^{(2)}&=v_j^{(1)}+s_jf(u_j^{(1)},\tau_j+\ell_j),\\
 u_j^{(2)}&=u_j^{(1)}+s_jf(v_j^{(2)},\tau_j+\ell_j).
\end{align*}
The independent root $\mu_j^*$ satisfies
\[
 \boxed{
 R_j(z^{[j-1]},\mu_j)
 =u_j^{(2)}-v_j^{(2)}+2\mu_j=0.
 }
\]
At the exact root,
\[
 z^{[j]}=u_j^{(2)}+\mu_j^*=v_j^{(2)}-\mu_j^*.
\]
The finite-tolerance implementation returns the neutral mean of these corrected
copies.  The analytic formulas below evaluate the ideal-root tangent at the
converged stages; a finite-difference derivative of the emitted map also
contains the stopping decision and the residual-level copy mismatch.

\section{Small GC2D matrices at the traversed stages}

For
\[
 f(z,t)=J_0\nabla_z\phi(z,t),
 \qquad
 J_0=\begin{pmatrix}0&-1\\1&0\end{pmatrix},
\]
introduce the signed small matrix
\[
 \boxed{
 w_j(z,t):=s_jD_zf(z,t)
 =s_j\begin{pmatrix}
 -\phi_{XY}&-\phi_{YY}\\
 \phi_{XX}&\phi_{XY}
 \end{pmatrix}_{(z,t)}.
 }
\]
The four matrices for substep $j$ are
\begin{align*}
 w_{j,1}&=w_j(v_j^{(0)},\tau_j),&
 w_{j,2}&=w_j(u_j^{(1)},\tau_j),\\
 w_{j,3}&=w_j(u_j^{(1)},\tau_j+\ell_j),&
 w_{j,4}&=w_j(v_j^{(2)},\tau_j+\ell_j),
\end{align*}
with
\[
 \Sigma_j=w_{j,2}+w_{j,3}.
\]
Because $s_2<0$, all four middle-substep matrices include that sign.  Replacing
them by absolute-step matrices differentiates a different map.

\section{Four duplicated-space stage Jacobians}

Each elementary shear has a $4\times4$ Jacobian containing one small block:
\[
 \boxed{
 \mathcal W_{j,1}=\begin{pmatrix}\Id_2&w_{j,1}\\0&\Id_2\end{pmatrix},
 \quad
 \mathcal W_{j,2}=\begin{pmatrix}\Id_2&0\\w_{j,2}&\Id_2\end{pmatrix},
 }
\]
\[
 \boxed{
 \mathcal W_{j,3}=\begin{pmatrix}\Id_2&0\\w_{j,3}&\Id_2\end{pmatrix},
 \quad
 \mathcal W_{j,4}=\begin{pmatrix}\Id_2&w_{j,4}\\0&\Id_2\end{pmatrix}.
 }
\]
The exact Jacobian of the unprojected signed ABBA map is the ordered product
\[
 \boxed{
 D\Phi^{\mathrm{ABBA}}_{\ell_j,\tau_j}
 =\mathcal W_{j,4}\mathcal W_{j,3}
  \mathcal W_{j,2}\mathcal W_{j,1}.
 }
\]
Writing it as
\[
 D\Phi^{\mathrm{ABBA}}_{\ell_j,\tau_j}
 =\begin{pmatrix}A_j&B_j\\C_j&D_j\end{pmatrix},
\]
the blocks are
\begin{align*}
 A_j&=\Id_2+w_{j,4}\Sigma_j,\\
 B_j&=w_{j,1}+w_{j,4}+w_{j,4}\Sigma_jw_{j,1},\\
 C_j&=\Sigma_j,\\
 D_j&=\Id_2+\Sigma_jw_{j,1}.
\end{align*}

\section{Local physical Jacobian}

Form the four $2\times2$ matrices
\begin{align*}
 K_j={}&4\Id_2-(w_{j,1}+w_{j,2}+w_{j,3}+w_{j,4})\\
       &+(w_{j,4}\Sigma_j+\Sigma_jw_{j,1})
        -w_{j,4}\Sigma_jw_{j,1},\\[1mm]
 L_j={}&w_{j,1}-w_{j,2}-w_{j,3}+w_{j,4}\\
       &+(w_{j,4}\Sigma_j-\Sigma_jw_{j,1})
        +w_{j,4}\Sigma_jw_{j,1},\\[1mm]
 P_j={}&\Id_2+w_{j,1}+w_{j,4}
          +w_{j,4}\Sigma_j+w_{j,4}\Sigma_jw_{j,1},\\[1mm]
 Q_j={}&2\Id_2-w_{j,1}-w_{j,4}
          +w_{j,4}\Sigma_j-w_{j,4}\Sigma_jw_{j,1}.
\end{align*}
Here
\[
 K_j=D_{\mu_j}R_j,
 \qquad
 L_j=D_{z^{[j-1]}}R_j.
\]
The implicit-function theorem gives
\[
 D\mu_j^*=-K_j^{-1}L_j,
\]
and the local physical Jacobian is
\[
 \boxed{
 J_j:=D\Psi^{[2]}_{\ell_j,\tau_j}(z^{[j-1]})
 =P_j-Q_jK_j^{-1}L_j.
 }
\]
The matrix $K_j$ is the reduced Newton matrix, not the physical Jacobian.  To
construct $J_j$ without an inverse, solve
\[
 K_jT_j=L_j,
 \qquad
 J_j=P_j-Q_jT_j.
\]

\section{Applying the triple-jump Jacobian}

The complete outer Jacobian is
\[
 \boxed{
 \mathcal J_n
 :=D\Psi^{[4]}_{h,t_n}(z_n)=J_3J_2J_1.
 }
\]
If only one tangent vector is needed, no $\mathcal J_n$ need be formed.  Start
with $\eta_0=\xi_n$ and repeat
\[
 b_j=L_j\eta_{j-1},
 \qquad
 K_jy_j=b_j,
 \qquad
 \eta_j=P_j\eta_{j-1}-Q_jy_j,
 \qquad j=1,2,3.
\]
Then
\[
 \boxed{\xi_{n+1}=\eta_3.}
\]
This procedure uses three $2\times2$ linear solves per particle but introduces
no additional nonlinear solve.

An equivalent local construction differentiates the four vector-field
increments.  With $M_j=D\mu_j^*$,
\begin{align*}
 Dk_{j,1}&=\frac{w_{j,1}}{s_j}(\Id_2-M_j),\\
 Du_j^{(1)}&=\Id_2+M_j+s_jDk_{j,1},\\
 Dk_{j,2}&=\frac{w_{j,2}}{s_j}Du_j^{(1)},\\
 Dk_{j,3}&=\frac{w_{j,3}}{s_j}Du_j^{(1)},\\
 Dv_j^{(2)}&=\Id_2-M_j+s_j(Dk_{j,2}+Dk_{j,3}),\\
 Dk_{j,4}&=\frac{w_{j,4}}{s_j}Dv_j^{(2)},
\end{align*}
and
\[
 J_j=\Id_2+\frac{\ell_j}{4}
 (Dk_{j,1}+Dk_{j,2}+Dk_{j,3}+Dk_{j,4}).
\]
This recurrence is useful as an independent analytic cross-check of the
implicit-function factorization.

\section{Observer-to-formula correspondence}

One accepted \texttt{ABBA4ImplicitIntegrationStep} contains the outer state,
time interval, coefficients, aggregate solver data, and an ordered
\texttt{substeps} tuple.  Each substep record supplies:
\begin{itemize}
 \item signed \texttt{start\_time} and \texttt{duration};
 \item \texttt{state\_before} and \texttt{state\_after};
 \item the converged multiplier;
 \item $u_j^{(0)},v_j^{(0)},u_j^{(1)},v_j^{(2)},u_j^{(2)}$;
 \item formulation, solver, iteration, residual, and tolerance metadata; and
 \item the dynamics object needed to evaluate the four spatial Jacobians.
\end{itemize}

The three snapshots are continuous:
\[
 z^{[0]}=z_n,
 \qquad
 z^{[j]}_{\mathrm{after}}=z^{[j+1]}_{\mathrm{before}},
 \qquad
 z^{[3]}=z_{n+1}.
\]
Shadow compositions used only to produce off-grid output samples do not emit
observation records and must not enter an accumulated tangent diagnostic.

\section{One-step diagnostic workflow}

For each observed outer step:
\begin{enumerate}
 \item Verify that exactly three substep records are present and that their
       coefficients, signed times, and boundary states are continuous.
 \item For each $j$, evaluate $w_{j,1},\ldots,w_{j,4}$ at the stored stage
       points and signed endpoint times.
 \item Form $K_j,L_j,P_j,Q_j$ in the stated order and solve the required local
       $2\times2$ systems.
 \item Propagate the tangent chronologically through $J_1$, $J_2$, and $J_3$,
       or construct $\mathcal J_n=J_3J_2J_1$.
 \item Compare the implicit-function and stage-increment factorizations when
       diagnosing analytic derivatives.
 \item Optionally compare with centered finite differences of the complete
       stored \texttt{map\_state}, choosing a perturbation above the nonlinear
       residual floor.
 \item Record local symplecticity and determinant defects without modifying the
       accepted physical trajectory.
\end{enumerate}

\section{Accumulated sensitivity}

To differentiate the numerical solution with respect to its initial state,
propagate $M_n\in\R^{2\times2}$:
\[
 M_0=\Id_2,
 \qquad
 \boxed{M_{n+1}=\mathcal J_nM_n.}
\]
After $N$ outer steps,
\[
 M_N=\mathcal J_{N-1}\cdots\mathcal J_1\mathcal J_0
 =\frac{\partial z_N}{\partial z_0}.
\]
The newest Jacobian multiplies from the left.  The update can be implemented
by applying the three local factorizations successively to the two columns of
$M_n$.

\section{Symplecticity diagnostics}

At exact projected roots, every local physical tangent is symplectic:
\[
 J_j^TJ_0J_j=J_0.
\]
Therefore,
\[
 \mathcal J_n^TJ_0\mathcal J_n=J_0,
 \qquad
 M_N^TJ_0M_N=J_0.
\]
In two dimensions this is equivalent to unit determinant.  Useful monitored
defects are
\[
 \varepsilon_{\mathrm{sp},n}
 =\norm{\mathcal J_n^TJ_0\mathcal J_n-J_0},
 \qquad
 \varepsilon_{\det,n}=\lvert\det\mathcal J_n-1\rvert.
\]
Their numerical size includes floating-point error, derivative accuracy, and
the residuals of all three nonlinear solves.  A small outer defect can mask
opposing local defects, so diagnostics should retain both the three $J_j$
values and their product when investigating a failure.

\section{Nearby-trajectory interpretation}

Let
\[
 z_n^\varepsilon=z_n+\varepsilon\xi_n.
\]
Differentiability of the complete implicit map gives
\[
 \Psi^{[4]}_{h,t_n}(z_n^\varepsilon)
 =\Psi^{[4]}_{h,t_n}(z_n)
 +\varepsilon\mathcal J_n\xi_n+\mathcal O(\varepsilon^2).
\]
Thus the integrator produces the reference trajectory while the optional
diagnostic produces the linear evolution of an infinitesimal separation.  The
negative middle duration changes the intermediate tangent, but the three
chronological applications still approximate the final nearby-trajectory
difference.

\section{Implementation-ready summary}

For every outer step, perform
\[
 \boxed{
 \begin{aligned}
 z_{n+1}&=\Psi^{[4]}_{h,t_n}(z_n)
 &&\text{with three independent projected solves},\\[1mm]
 \xi_{n+1}&=J_3J_2J_1\xi_n
 &&\text{with three local ideal-root tangents}.
 \end{aligned}
 }
\]
For each substep, compute
\[
 K_jy_j=L_j\eta_{j-1},
 \qquad
 \eta_j=P_j\eta_{j-1}-Q_jy_j.
\]
The physical state is never obtained by multiplying its previous value by a
Jacobian.  The tangent is a read-only diagnostic derived from the already
accepted multipliers and stage snapshots.

\subsection*{Minimum retained data}

A reproducible outer-step record must preserve the ordered coefficient vector,
the three signed start times and durations, the three accepted boundary states,
and each substep's converged multiplier and stage snapshots.  Keeping only the
aggregate iteration count or the maximum multiplier norm is insufficient to
reconstruct $J_1$, $J_2$, and $J_3$.

\subsection*{Interpreting discrepancies}

When two Jacobian diagnostics disagree, check the following causes in order:
\begin{itemize}
 \item a lost negative sign in $s_2=\delta h/2$ or a stage evaluated at the
       outer time instead of its signed local endpoint;
 \item reordered matrix products, especially $w_{j,4}\Sigma_jw_{j,1}$ or the
       outer product $J_3J_2J_1$;
 \item comparison of an ideal-root analytic tangent with a finite-difference
       perturbation below the nonlinear residual floor;
 \item an ill-conditioned $K_j$, which amplifies residual, derivative, and
       round-off errors; or
 \item accidental inclusion of an unobserved shadow composition in the
       accumulated tangent history.
\end{itemize}

These checks localize the defect to one signed solve before the complete outer
composition is assessed.

\end{document}
